\documentclass[11pt]{article}
\usepackage[utf8]{inputenc}
\usepackage[margin=1in, headheight=15pt]{geometry}
\usepackage{amsmath, amssymb, amsfonts}
\usepackage{fancyhdr}
\usepackage{enumitem}

% Setup page headers and footers to match the uploaded PDF style
\pagestyle{fancy}
\fancyhf{}
\lhead{CSE 400: Fundamentals of Probability in Computing}
\rhead{Lecture Scribe}
\cfoot{\thepage}

\begin{document}

% Title Block formatted similar to the PDF's opening lines
\begin{flushleft}
\renewcommand{\arraystretch}{1.2}
\begin{tabular}{@{}ll}
\textbf{Course:} & CSE400 – Fundamentals of Probability in Computing \\
\textbf{Lecture Topic:} & Marginal PMF, Correlation, Covariance, Joint Gaussian Random Variables \\
\textbf{Date:} & March 12, 2026 \\
\textbf{Scribe:} & Het Patel \\
\end{tabular}
\end{flushleft}

\vspace{0.3cm}

\section{Opening Context (Real-Life Connection)}
(Explicitly mentioned in lecture slides)

Joint Gaussian models are widely used in:
\begin{itemize}
    \item Signal processing
    \item Communication systems
    \item Modeling noisy real-world measurements
\end{itemize}

\section{Marginal PMF}
\subsection{Definition}
Definition (as stated in lecture):\\
For discrete random variables $X$ and $Y$, the joint PMF is:
\[ p_{XY}(x,y) = P(X = x, Y = y) \]

The marginal PMFs are:
\[ p_X(x) = \sum_{y} p_{XY}(x,y) \]
\[ p_Y(y) = \sum_{x} p_{XY}(x,y) \]

\textbf{Key Idea (from lecture):}
The marginal PMF of one variable is obtained by summing the joint PMF over all possible values of the other variable.

\subsection{Notation}
\begin{itemize}
    \item $p_{XY}(x,y)$: Joint PMF
    \item $p_X(x), p_Y(y)$: Marginal PMFs
    \item Summation over other variable removes dependency
\end{itemize}

\subsection{Marginal PMF from Joint PMF Table}
(From table shown in lecture slides)
\begin{itemize}
    \item Row sums $\rightarrow p_X(x)$
    \item Column sums $\rightarrow p_Y(y)$
\end{itemize}
\[ p_X(x_i) = \sum_j p_{ij}, \quad p_Y(y_j) = \sum_i p_{ij} \]
\[ \sum_i p_X(x_i) = 1, \quad \sum_j p_Y(y_j) = 1 \]

\subsection{Worked Example (from lecture slides)}
Given Joint PMF:
\begin{center}
\begin{tabular}{c|cc}
      & $y=0$ & $y=1$ \\
\hline
$x=0$ & $\frac{1}{8}$ & $\frac{1}{4}$ \\[6pt]
$x=1$ & $\frac{3}{8}$ & $\frac{1}{4}$ \\
\end{tabular}
\end{center}

\textbf{Step 1: Compute $p_X(x)$}
\[ p_X(0) = \frac{1}{8} + \frac{1}{4} = \frac{3}{8} \]
\[ p_X(1) = \frac{3}{8} + \frac{1}{4} = \frac{5}{8} \]

\textbf{Step 2: Compute $p_Y(y)$}
\[ p_Y(0) = \frac{1}{8} + \frac{3}{8} = \frac{1}{2} \]
\[ p_Y(1) = \frac{1}{4} + \frac{1}{4} = \frac{1}{2} \]

\textbf{Final Answer:}
\[
p_X(x) =
\begin{cases}
\frac{3}{8}, & x=0 \\
\frac{5}{8}, & x=1
\end{cases}
\quad
p_Y(y) =
\begin{cases}
\frac{1}{2}, & y=0 \\
\frac{1}{2}, & y=1
\end{cases}
\]

\section{Correlation (Raw Measure)}
\subsection{Definition (from lecture)}
\[ R_{XY} = E[XY] \]

\subsection{Interpretation}
\begin{itemize}
    \item Measures average value of product $XY$
    \item Computed about origin (not centered)
    \item Depends on: Mean, Spread, Dependence
\end{itemize}

\subsection{Important Observation}
$R_{XY} = 0 \Rightarrow$ orthogonal about origin \\
Not a clean dependence measure \\

\textbf{Transition (as emphasized in lecture):}
To remove mean effects $\rightarrow$ leads to Covariance

\section{Covariance}
\subsection{Definition}
\[ \text{Cov}(X,Y) = E[(X - \mu_X)(Y - \mu_Y)] \]
where:
\[ \mu_X = E[X], \quad \mu_Y = E[Y] \]

\subsection{Algebraic Form (Derived in Lecture)}
Step-by-step expansion:
\begin{align*}
\text{Cov}(X,Y) &= E[XY - X\mu_Y - Y\mu_X + \mu_X\mu_Y] \\
&= E[XY] - \mu_Y E[X] - \mu_X E[Y] + \mu_X\mu_Y \\
&= E[XY] - \mu_X \mu_Y
\end{align*}

\subsection{Final Formula}
\[ \text{Cov}(X,Y) = E[XY] - E[X]E[Y] \]

\subsection{Interpretation (from slides)}
\begin{itemize}
    \item Positive covariance $\rightarrow$ variables increase together
    \item Negative covariance $\rightarrow$ one increases, other decreases
    \item Zero covariance $\rightarrow$ no linear relationship
\end{itemize}
\textbf{Warning:} Magnitude depends on units $\rightarrow$ cannot measure strength

\section{Properties of Covariance}
\begin{itemize}
    \item \textbf{Symmetry:} $\text{Cov}(X,Y) = \text{Cov}(Y,X)$
    \item \textbf{Variance as Special Case:} $\text{Cov}(X,X) = \text{Var}(X)$
    \item \textbf{Linearity:} $\text{Cov}(aX+b, cY+d) = ac\,\text{Cov}(X,Y)$
    \item \textbf{Independence $\Rightarrow$ Zero Covariance:} $X \perp Y \Rightarrow \text{Cov}(X,Y) = 0$
\end{itemize}
\textbf{Important Warning (from lecture):} Zero covariance $\neq$ independence

\section{Worked Example 1 (from lecture)}
\textbf{Problem:}
If $Y = -2X + 3$, find $\text{Cov}(X,Y)$

\textbf{Solution:}
\[ \text{Cov}(X,Y) = E[XY] - E[X]E[Y] \]

\textbf{Step 1: Compute $XY$}
\[ XY = X(-2X + 3) = -2X^2 + 3X \]

\textbf{Step 2: Compute expectations}
\[ E[XY] = E[-2X^2 + 3X] = -2E[X^2] + 3E[X] \]
\[ E[Y] = E[-2X + 3] = -2E[X] + 3 \]

\textbf{Step 3: Substitute}
\begin{align*}
\text{Cov}(X,Y) &= (-2E[X^2] + 3E[X]) - E[X](-2E[X] + 3) \\
&= -2E[X^2] + 3E[X] + 2(E[X])^2 - 3E[X] \\
&= -2E[X^2] + 2(E[X])^2
\end{align*}

\textbf{Step 4: Use variance identity}
\[ \text{Var}(X) = E[X^2] - (E[X])^2 \]
\[ \Rightarrow \text{Cov}(X,Y) = -2\text{Var}(X) \]

\section{Pearson’s Correlation Coefficient}
\subsection{Definition}
\[ \rho_{XY} = \frac{\text{Cov}(X,Y)}{\sqrt{\text{Var}(X)\text{Var}(Y)}} = \frac{\text{Cov}(X,Y)}{\sigma_X \sigma_Y} \]

\subsection{Properties}
\begin{itemize}
    \item $-1 \le \rho_{XY} \le 1$
    \item Unit-free
    \item Measures strength + direction
\end{itemize}

\subsection{Interpretation}
\begin{itemize}
    \item $\rho > 0$: Positive linear trend
    \item $\rho < 0$: Negative linear trend
    \item $\rho = 0$: No linear correlation
\end{itemize}

\section{Worked Example 2 (Continuous Case)}
\textbf{Given:}
\[
f_{XY}(x,y) =
\begin{cases}
\frac{xy}{96}, & 0 < x < 4,\ 1 < y < 5 \\
0, & \text{otherwise}
\end{cases}
\]

\textbf{Key Results (as derived in lecture slides)}
\[ E[X] = \frac{8}{3}, \quad E[Y] = \frac{31}{9} \]
\[ E[XY] = \frac{248}{27} \]
\[ E[2X+3Y] = \frac{47}{3} \]
\[ \text{Var}(X) = \frac{8}{9}, \quad \text{Var}(Y) = \frac{92}{81} \]

\textbf{Covariance}
\begin{align*}
\text{Cov}(X,Y) &= E[XY] - E[X]E[Y] \\
&= \frac{248}{27} - \left(\frac{8}{3} \cdot \frac{31}{9}\right) = 0
\end{align*}

\textbf{Correlation}
\[ \rho_{XY} = 0 \]

\section{Joint Gaussian Random Variables}
\subsection{Definition}
A pair $(X,Y)$ is jointly Gaussian if:
\begin{itemize}
    \item Joint PDF has bivariate Gaussian form
    \item Marginals are Gaussian
\end{itemize}

\subsection{Joint PDF (as given in lecture)}
\[
f_{XY}(x,y) =
\frac{1}{2\pi \sigma_X \sigma_Y \sqrt{1-\rho^2}}
\exp\left(
-\frac{1}{2(1-\rho^2)}
\left[
\left(\frac{x-\mu_X}{\sigma_X}\right)^2
+
\left(\frac{y-\mu_Y}{\sigma_Y}\right)^2
-
2\rho \left(\frac{x-\mu_X}{\sigma_X}\right)\left(\frac{y-\mu_Y}{\sigma_Y}\right)
\right]
\right)
\]

\subsection{Parameters}
\begin{itemize}
    \item $\mu_X, \mu_Y$: Means
    \item $\sigma_X, \sigma_Y$: Standard deviations
    \item $\rho$: Correlation coefficient
\end{itemize}

\subsection{Interpretation (from lecture example)}
If $\rho > 0$: variables increase together \\
Example given: CO$_2$ concentration vs temperature \\
Higher CO$_2 \rightarrow$ higher temperature trend

\section{Connection to Previous Material}
\begin{itemize}
    \item \textbf{Uses:} Expectation $E[X]$, Variance definition
    \item \textbf{Extends:} Single RV $\rightarrow$ Joint RV analysis
    \item \textbf{Leads to:} Multivariate distributions
\end{itemize}

\section{Summary \& Key Takeaways}
\begin{itemize}
    \item Marginal PMF = sum over other variable
    \item $E[XY]$ is raw correlation (not centered)
    \item Covariance removes mean effect
    \item Covariance: sign $\rightarrow$ direction, not strength
    \item Pearson coefficient: normalized measure, gives strength + direction
    \item Zero covariance $\neq$ independence
    \item Joint Gaussian: fundamental continuous model
\end{itemize}

\newpage

\section*{APPENDIX A: Quick Reference Sheet}

\textbf{Formulas}
\vspace{0.3cm}

\noindent
\renewcommand{\arraystretch}{1.5}
\begin{tabular}{|l|l|}
\hline
\textbf{Concept} & \textbf{Formula} \\
\hline
Marginal PMF & $p_X(x) = \sum_y p_{XY}(x,y)$ \\
\hline
Correlation (raw) & $E[XY]$ \\
\hline
Covariance & $E[XY] - E[X]E[Y]$ \\
\hline
Pearson Coefficient & $\rho = \frac{\text{Cov}}{\sigma_X \sigma_Y}$ \\
\hline
Joint Gaussian PDF & (as given in section 9.2) \\
\hline
\end{tabular}

\vspace{0.5cm}
\textbf{Key Conditions}
\begin{itemize}
    \item Marginalization $\rightarrow$ sum over other variable
    \item Covariance uses centered variables
    \item Independence $\Rightarrow$ covariance = 0
    \item \textbf{Warning:} Converse not true
\end{itemize}

\vspace{0.3cm}
\textbf{Buzz Words (from lecture)}
\begin{itemize}
    \item ``Summing out'' $\rightarrow$ marginalization
    \item ``Centered variables'' $\rightarrow$ covariance
    \item ``Normalized'' $\rightarrow$ Pearson coefficient
    \item ``Linear relationship'' $\rightarrow$ correlation
\end{itemize}

\vspace{0.3cm}
\textbf{Space for Intuition (Not in Lecture)}
\begin{itemize}
    \item Marginal PMF = ``ignore the other variable''
    \item Covariance = ``do deviations move together?''
    \item Pearson = ``scaled covariance''
\end{itemize}

\end{document}